\chapter{Riemann 流形}\label{chap:riemannian-manifolds}
\begingroup
\renewcommand{\thefootnote}{}
\footnotetext{\textcopyright\ Springer International Publishing Switzerland 2014.
  L. Godinho and J. Natário, \emph{An Introduction to Riemannian Geometry},
  Universitext. DOI \href{https://doi.org/10.1007/978-3-319-08666-8_3}{10.1007/978-3-319-08666-8\_3}.}
\endgroup

$\mathbb{R}^n$ 的度量性质 (距离和角度) 由其典范的 Cartesian 坐标确定. 然而, 一般的可微流形上没有这样一套优先选取的坐标. 要定义距离和角度, 必须选定一个特殊的 2-张量场, 加入更多结构. 这个张量场称为 \textbf{Riemann 度量} (这与必须先选取一个体积形式才能确定体积概念十分相似). Gauss, Bolyai 和 Lobachevsky 在约 1830 年发现非 Euclidean 几何之后, Riemann 在 1854 年的任教资格演讲\textquotedblleft 论作为几何学基础的假设\textquotedblright 中提出这一思想 (实际上, Riemann 的演讲题目正是由 Gauss 提议的). 这一概念被证明极富成效, 其成果之一是促成了 Einstein 广义相对论的发展.

本章开始研究 Riemann 几何. 第~\RGref{sec:riemannian-metrics}~节将 Riemann 度量引入为给每个切空间确定一个内积的张量场. 这自然引出许多概念, 如向量 (或曲线) 的\textbf{长度}, 两个向量之间的\textbf{夹角}, \textbf{Riemann 体积形式} (赋予任意标准正交基以单位体积) 和函数的\textbf{梯度}.

第~\RGref{sec:affine-connections}~节讨论向量场的微分. 这一概念也需要引入某种额外的结构, 称为\textbf{仿射联络}, 因为可微流形上的向量场没有可优先选取的 Cartesian 分量可供求导. 仿射联络给出沿曲线的向量\textbf{平行性}概念, 进而定义\textbf{测地线}, 即切向量平行的曲线. Riemann 流形带有一种特殊的仿射联络, 称为 \textbf{Levi-Civita 联络} (第~\RGref{sec:levi-civita-connection}~节), 其测地线具有距离最小化性质 (第~\RGref{sec:minimizing-properties-of-geodesics}~节). 这与两点之间最短的路线不转弯这一直观想法一致.

最后, 第~\RGref{sec:hopf-rinow-theorem}~节证明 \textbf{Hopf-Rinow 定理}, 它把 Riemann 流形作为度量空间的性质与其测地线的性质联系起来. 该定理完全刻画了\textbf{完备} Riemann 流形这一重要类别.
